\chapter{The verdict, and what would have to change}\label{ch:hslverdict}
\def\tbrk{\hfil\penalty50\hfilneg}% local: a line may break before a provenance tag

\begin{onething}
HyperSL is a sound model of what spectra look like inside one scene, and on every one of this product's jobs a design
choice that serves that purpose (per-scene scaling, band masking within a pixel, a tokenizer over band index, smoothing
heads, cloud-free pretraining) works against the job; making it fit would mean retraining it with our choices, which is
what Indradhanu already is.
\end{onething}

This chapter puts \cref{ch:hslwhat,ch:hslnorm,ch:hsldetect,ch:hslunmix,ch:hslcloud} side by side. It adds no new
measurement; every number has its source in the chapter that measured it.

\begin{analogy}
A superb road bicycle offered to a team that delivers parcels across a ploughed field. Everything that makes it superb
(thin tyres, no luggage rack, a riding position for speed) is a deliberate choice for a road, and each one is a problem in
a field. You could fit fat tyres, a rack and a new frame; at that point you have built a different bicycle, and the team
already owns one.
\end{analogy}

%======================================================================
\section{The decision table}\label{sec:hslverdict:table}

\begin{center}\small
\begin{longtable}{@{}>{\raggedright\arraybackslash}p{30mm}>{\raggedright\arraybackslash}p{40mm}>{\raggedright\arraybackslash}p{36mm}>{\raggedright\arraybackslash}p{26mm}@{}}
\toprule
requirement & HyperSL & this product & evidence \\
\midrule\endhead
the same spectrum has the same input in every scene & no: per-band min-max over the scene &
\code{WavelengthStats}, one fixed curve & \cref{we:hslnorm:s1,we:hslnorm:s2} \\
a library spectrum has an input & none; the scene frame clipped 41\% of ECOSTRESS & the same curve, the same stem &
\cref{sec:hslnorm:toy} \\
anomalies matched to man-made library entries & below SAM on 3 of 4 scenes; one leaf took 59.6\% of set2 &
Simplex ACE (adopted), the namer (in training) & \cref{tab:hslnorm:m1} \\
band lists with removed windows & convolution over index; a 5-band window spans 217\,nm & 40\,nm wavelength bins &
\cref{tab:hslnorm:spans} \\
a pixel must not help rebuild itself & rebuilds a pixel's bands from its own bands & hides whole pixels, rebuilds from
neighbours & \cref{we:hsldetect:self} \\
detection on the 909-object bench & own residual: 130 objects, pAUC 0.018; with our scorer: 350, pAUC 0.083, below PCA-RX's 0.128 with no
model & FTH3 ep.\ 5: pAUC 0.189, 449 objects & \cref{tab:hsldetect:main} \\
sub-pixel targets kept & 0.25\,m target: object AUC 0.68 & ACE 0.97 & \cref{tab:hslunmix:hysu} \\
blind unmixing on exact truth & worse than VCA + FCLS on every set & VCA + FCLS & \cref{tab:hslunmix:msad} \\
thin cloud found & 28.9\% of EnMAP's cirrus & 67.9\% from raw bands; Indradhanu's head & \cref{we:hslcloud:cirrus} \\
\bottomrule
\end{longtable}
\end{center}

\subsection{Where HyperSL did better}

A verdict that lists only losses is not a fair one. In the same experiments HyperSL was ahead on five measures:
\begin{itemize}
\item its embedding groups a spectral library slightly better than raw angle (5-neighbour same-class share 0.965 against
0.949 on ECOSTRESS, \cref{sec:hslnorm:measured});
\item the cloud classifier on its features found more thick cloud (87.1\% against 80.4\%) and made fewer false calls on the
high plateau (4.34\% against 6.62\%, \cref{ch:hslcloud});
\item as a second-stage filter on Indradhanu's candidates, its embeddings removed more false alarms at 90\% object recall
when the cut was set with hindsight (\cref{sec:hsldetect:filter});
\item on Samson, Jasper and Urban it agreed with the FCLS references better than VCA + FCLS (\cref{tab:hslunmix:msad});
\item it reproduces its paper's clustering and fine-tuning claims (\cref{sec:hslwhat:repro}).
\end{itemize}
None of these is one of the product's jobs as the product does them, and the third did not survive an honest threshold.

%======================================================================
\section{What it costs}\label{sec:hslverdict:cost}

\begin{center}\small
\begin{tabular}{@{}>{\raggedright\arraybackslash}p{24mm}>{\raggedright\arraybackslash}p{110mm}@{}}
\toprule
cost & what the records show \\
\midrule
compute & HyperSL needs a GPU: 1,773\,s of laptop-GPU time for the four bench scenes, 438 microseconds per pixel at
4 draws (\cref{sec:hsldetect:run}); Indradhanu Heavy runs at 454\,s per scene on 16 CPU threads in the product worker. The unmixing variant took 5--13 minutes per scene on a 96\,GB GPU, against seconds for VCA + FCLS. \\
licence & no LICENSE file on GitHub; the Hugging Face card says MIT. Unresolved for a product. \\
data & the $\sim$265\,GB pretraining corpus is on Baidu Netdisk only. We cannot audit what it learned from, or retrain it
with our scaling, without it. \\
paper against code & seven surprises found only by reading the code: base 100 against 10,000 in the position code; ``layer-by-layer normalization'' that
is per-scene min-max; 300 against 200 million spectra; a 439-band airborne set not in \file{train.py}; the
``endmember-free'' code that starts from VCA; downstream scripts that never load the checkpoint; a final encoder layer with
no skip connection. \\
\bottomrule
\end{tabular}
\end{center}

%======================================================================
\section{What would have to change}\label{sec:hslverdict:change}

Each obstacle has a fix, and each fix is a change to how the model is \emph{trained}, not to how it is called.
\begin{enumerate}
\item \textbf{Scaling.} Train on reflectance with a fixed per-wavelength normalisation, so a pixel's input is a property of
the pixel and a library spectrum has one (\cref{ch:hslnorm}).
\item \textbf{Tokenizer.} Make the first layer read wavelengths, not band positions, so a removed window is a gap rather
than a neighbour (\cref{sec:hslnorm:conv}).
\item \textbf{Masking and context.} Hide whole pixels and give the model their neighbours, so it cannot complete an anomaly
from itself (\cref{ch:hsldetect}).
\item \textbf{Data.} Keep clouds in pretraining and label them (\cref{ch:hslcloud}).
\item \textbf{Heads.} Judge any per-pixel head at fills of 0.1--0.25, and do not smooth abundances across pixels
(\cref{ch:hslunmix}).
\item \textbf{Access.} Publish the training corpus somewhere it can be read, and a licence in the code repository.
\end{enumerate}
Items 1--4 together describe a wavelength-aware spatial masked autoencoder pretrained on reflectance with fixed
statistics: Indradhanu on its SpectralEarth stem (\cref{ch:specdecoder,ch:pretrain}). The honest summary is not that
HyperSL is bad. It is that \textbf{the changes it would need are the design we already chose, for the same reasons.}

\subsection{What we took from it}

\begin{itemize}
\item \textbf{Wavelength as a coordinate.} HyperSL's position code accepts any band list; our stem goes further and builds
its bins from wavelength, but the idea is shared.
\item \textbf{A discipline for outside models.} Check the checkpoint's hash, read the code before the paper, reproduce a
published number before measuring anything new, and keep the outputs. Every claim in this part could be checked because
those steps were followed, and the one that could not (the 28\% of \cref{sec:hslnorm:measured}) is the one whose output
was not kept.
\end{itemize}

\begin{drift}
Every number in the decision table belongs to a dated run (2026-09-23 to 2026-10-31) and to HyperSL's \code{10_base_mask95}
release. A new release (different scaling, spatial context, a licence) would reopen each row.
\end{drift}

%======================================================================
\begin{confusions}
\confusion{The part shows HyperSL is a bad model.}{It shows a good model for one-scene land-cover description, measured on
jobs it was not built for.}
\confusion{Fine-tuning would fix it.}{Fine-tuning keeps the input scaling, the band masking and the tokenizer. The obstacles
are in what the model was trained to do, so they need a retrain.}
\confusion{Its false-alarm filter result means it helps detection.}{The gain needed a threshold chosen with the test scene's
labels; with a threshold set on other scenes it lost objects (\cref{sec:hsldetect:filter}).}
\confusion{The licence is fine because Hugging Face says MIT.}{The code repository has no licence file. For a product, that
needs settling with the authors.}
\end{confusions}

\begin{checkyourself}
\item\label{ex:hslverdict:1} For each of the four jobs (matching, detection, unmixing, cloud), name the single design choice of
  HyperSL most responsible for its result, and the chapter that shows it.
\item\label{ex:hslverdict:2} Which of the six changes in \cref{sec:hslverdict:change} could be made without retraining? Explain.
\item\label{ex:hslverdict:3} From \cref{tab:hsldetect:main}, how many more objects does FTH3 epoch 5 catch in the top 1\% than HyperSL RM,
  both with cloud excluded? What share of the 909 is that, and which bootstrap row agrees?
\item\label{ex:hslverdict:4} List the five measures on which HyperSL was ahead, and say for each why it does not change the
  decision.
\end{checkyourself}
